\documentclass[11pt,a4paper]{article}
\usepackage[T1]{fontenc}
\usepackage[utf8]{inputenc}
\usepackage{amsmath,amssymb,amsthm}
\usepackage[margin=1in]{geometry}
\usepackage[colorlinks=true,linkcolor=blue,urlcolor=blue]{hyperref}
\theoremstyle{plain}
\newtheorem{theorem}{Theorem}
\newtheorem{lemma}[theorem]{Lemma}
\newtheorem{proposition}[theorem]{Proposition}
\newtheorem{corollary}[theorem]{Corollary}
\theoremstyle{definition}
\newtheorem{remark}[theorem]{Remark}

\title{The empirical recurrence for OEIS A206006, proved by transfer matrix}
\author{Adrian Perez Fontelles\\ \small Independent researcher}
\date{2 September 2026}
\begin{document}
\maketitle

\begin{abstract}
OEIS A206006 counts arrays over $\{0,\dots,2\}$ under a condition on the number of clockwise edge increases
in every $2\times2$ subblock, and carries an empirical recurrence of order $38$ contributed by
R.~H.~Hardin. It is true, and it is decidable rather than empirical. The condition is local to
a bounded window of consecutive rows, so the count is a walk count on $S=243$ states and is
$C$-finite; whether the conjectured recurrence annihilates it is then settled by a finite exact
computation, with the Cayley--Hamilton theorem bounding the work.
\end{abstract}

\noindent\small 2020 Mathematics Subject Classification. 05A15, 05B45, 15B36.\normalsize

\section{The conjecture}

OEIS A206006 is ``Number of (n+1) X 5 0..2 arrays with the number of clockwise edge increases in every 2 X 2 subblock equal to one, and every 2 X 2 determinant nonzero.''. It has offset $1$ and begins
\[
722, 3706, 12068, 64224, 211824, 1159120, 3850160, 21563332,\ \dots
\]
The entry states, as a conjecture contributed by its author and never marked settled:
\begin{quote}
Empirical: a(n) = 8*a(n-1) +42*a(n-2) -482*a(n-3) -293*a(n-4) +11440*a(n-5) -11229*a(n-6) -137742*a(n-7) +261810*a(n-8) +894760*a(n-9) -2462973*a(n-10) -2981208*a(n-11) +12607691*a(n-12) +3180168*a(n-13) -38516874*a(n-14) +11294550*a(n-15) +72866395*a(n-16) -51223616*a(n-17) -84265347*a(n-18) +98230578*a(n-19) +50978982*a(n-20) -112745308*a(n-21) +3184853*a(n-22) +82656280*a(n-23) -34966370*a(n-24) -37445012*a(n-25) +32301744*a(n-26) +8331808*a(n-27) -15845182*a(n-28) +649616*a(n-29) +4536888*a(n-30) -914656*a(n-31) -719464*a(n-32) +226432*a(n-33) +51024*a(n-34) -24320*a(n-35) -128*a(n-36) +1024*a(n-37) -128*a(n-38) for n>40.
\end{quote}
As of the ``Last modified'' line on the live entry (Jul 17 2025 19:09:41, revision 7) this is still
recorded as empirical, and nothing on the entry records it as proved.

\section{What is being counted}

For a $2\times2$ subblock $\begin{pmatrix}p&q\\r&s\end{pmatrix}$ the
entry's quantity, the number of clockwise edge increases, counts how many of the four steps of the cycle $p\to q\to s\to r\to p$ go up in value. Write $\sigma$ for it. The subblocks of an
$(n+1)\times5$ array over $\{0,\dots,2\}$ form a grid of $n$ rows and $4$
columns and $\sigma$ labels that grid.

The reading is not a guess. Over $0..1$ the $2\times3$ arrays whose two subblocks have equal
clockwise counts number $40$; over $0..2$ the $2\times2$ arrays of clockwise count $2$ number
$30$; the $3\times2$ binary arrays whose two subblocks have unequal rightwards-and-downwards
counts number $40$. Those are the first terms the corresponding entries publish, and each was
computed by direct enumeration before any transfer matrix was built.

The condition is on each subblock separately: $\sigma$ must equal $1$ and, the entry adds, no subblock may be singular. Nothing
relates one subblock to another.

\section{The count is a walk count}

A subblock is built from two consecutive array rows and the condition looks
at one subblock at a time, so a single row is a state: rows $r$ and $s$ may follow one another
exactly when all $4$ subblocks they form pass. An $(n+1)$-row array is a walk of $n$ steps and
\[
a(n)\;=\;\iota^{\!\top}M^{\,n}\tau,\qquad\iota=\tau=(1,\dots,1)^{\!\top},
\]
on the $S=243$ rows.

At $n=1$ the array has two rows, so the grid is a single row of $4$
subblocks; the entry's first term is $722$, which the model
reproduces along with every later term. That is the cheapest check that the grid has been read
with the right shape.
In particular $a$ is $C$-finite. The alignment of the walk index with the entry's own $n$ was
checked against its published terms rather than assumed.

\section{The proof}

\begin{lemma}\label{lem:crit}
Let $q(t)=t^{38}-\sum_i c_it^{38-i}$ be the characteristic polynomial of the
conjectured recurrence. The residual $a(n)-\sum_ic_ia(n-i)$ equals
$\iota^{\!\top}M^{\,j}q(M)\tau$ for the corresponding walk index $j$, so the recurrence
holds from some point on if and only if $\iota^{\!\top}M^{j}q(M)\tau=0$ for all large $j$,
and it is enough to examine $j<S$.
\end{lemma}

\begin{proof}
Factor $M^{j}$ out of $M^{j+38}-\sum_ic_iM^{j+38-i}$. For the bound, $M^{S}$
is an integer combination of $I,M,\dots,M^{S-1}$ by Cayley--Hamilton, so the numbers
$u_j=\iota^{\!\top}M^{j}q(M)\tau$ obey the monic recurrence given by the characteristic
polynomial of $M$; $S$ consecutive zeros force every later one, and the last nonzero $u_j$
pins the threshold exactly.
\end{proof}

\begin{theorem}
\[
a(n)\;=\;8\,a(n-1) + 42\,a(n-2) - 482\,a(n-3) - 293\,a(n-4) + 11440\,a(n-5) - 11229\,a(n-6) - 137742\,a(n-7) + 261810\,a(n-8) + 894760\,a(n-9) - 2462973\,a(n-10) - 2981208\,a(n-11) + 12607691\,a(n-12) + 3180168\,a(n-13) - 38516874\,a(n-14) + 11294550\,a(n-15) + 72866395\,a(n-16) - 51223616\,a(n-17) - 84265347\,a(n-18) + 98230578\,a(n-19) + 50978982\,a(n-20) - 112745308\,a(n-21) + 3184853\,a(n-22) + 82656280\,a(n-23) - 34966370\,a(n-24) - 37445012\,a(n-25) + 32301744\,a(n-26) + 8331808\,a(n-27) - 15845182\,a(n-28) + 649616\,a(n-29) + 4536888\,a(n-30) - 914656\,a(n-31) - 719464\,a(n-32) + 226432\,a(n-33) + 51024\,a(n-34) - 24320\,a(n-35) - 128\,a(n-36) + 1024\,a(n-37) - 128\,a(n-38)
\]
for every $n>40$.
\end{theorem}

\begin{proof}
The vector $w=q(M)\tau$ was formed by $38$ matrix--vector products in exact integer
arithmetic and $\iota^{\!\top}M^{j}w$ evaluated for $j=0,1,\dots$, again exactly; those
integers vanish from the index corresponding to $n=40+1$ onwards and the run continued
until the working vector was identically zero, which settles every larger $j$ at once.
\end{proof}

\section{Verification}

Three checks, each independent of the algebra above.

First, the model reproduces all $19$ terms the entry publishes, exactly, in integer
arithmetic. This is what ties the model to the entry: the digraph is built by reading the
entry's English, and a misreading gives different counts.

Second, the conjectured recurrence was evaluated directly on those published terms, with no
matrices involved, and holds wherever the proved range and the data overlap.

Third, the annihilation test of Lemma~\ref{lem:crit} was carried out in exact integer
arithmetic on the whole vector rather than on a sampled prefix.

\begin{thebibliography}{9}
\bibitem{oeis} The OEIS Foundation, \emph{The On-Line Encyclopedia of Integer
Sequences}, \url{https://oeis.org}, 2026. Sequence A206006.
\bibitem{stanley} R.~P.~Stanley, \emph{Enumerative Combinatorics}, Volume 1, 2nd ed.,
Cambridge University Press, 2011. (Transfer-matrix method, Section 4.7.)
\bibitem{bona} M.~B\'ona, \emph{Combinatorics of Permutations}, 2nd ed., CRC
Press, 2012.
\bibitem{horn} R.~A.~Horn and C.~R.~Johnson, \emph{Matrix Analysis}, 2nd ed., Cambridge
University Press, 2013.
\end{thebibliography}

\end{document}
